% answers.tex —— 1958 年·云南卷 · 参考答案与详细解析片段
% 由 tools/generate_projects.py 依来源总账自 1958/1958_云南_数学真题_试题.md + 1958/1958_云南_数学真题_解析正文.md + 1958/1958_云南_数学真题_答案要点.md 生成，勿手改（改源 MD 后重跑）


\begin{gknotice}
\textbf{解析归属}：本卷题面出自固定原卷（仅题面，无答案页）；以下各题解析由本库独立推导，不声称代表官方发布版本。原卷 $\log$ 未标底数，解析按常用对数（底 $10$）书写；若底为 $a$，把 $10$ 换成 $a$ 即可。
\end{gknotice}

\begin{answer}
\gksolhead{1}{}

\textbf{题目}：

解方程 $\left(\log x^{\sqrt5}\right)^2+\log x^5-\dfrac54=0$。

\Needspace{4\baselineskip}
\textbf{解析}

\Needspace{5\baselineskip}
设 $t=\log x$（定义域要求 $x>0$）。由幂的对数性质，$\log x^{\sqrt5}=\sqrt5\,\log x=\sqrt5\,t$，$\log x^5=5\log x=5t$，原方程化为

\nopagebreak[4]
\[
\left(\sqrt5\,t\right)^2+5t-\frac54=0
\]

\Needspace{5\baselineskip}
即 $5t^2+5t-\dfrac54=0$，两边除以 $5$：

\nopagebreak[4]
\[
t^2+t-\frac14=0
\]

\Needspace{5\baselineskip}
解得

\nopagebreak[4]
\[
t=\frac{-1\pm\sqrt{1+1}}{2}=\frac{-1\pm\sqrt2}{2}
\]

\Needspace{5\baselineskip}
两个 $t$ 值都使 $x=10^t>0$，均在定义域内。故

\nopagebreak[4]
\[
x_1=10^{\frac{\sqrt2-1}{2}},\qquad x_2=10^{-\frac{\sqrt2+1}{2}}
\]

\gkanswer{答案：$x=10^{\frac{\sqrt2-1}{2}}$ 或 $x=10^{-\frac{\sqrt2+1}{2}}$}

\end{answer}
\begin{answer}
\gksolhead{2}{}

\textbf{题目}：

求 $(1+x)+(1+x)^2+\cdots+(1+x)^{10}$ 式中 $x^8$ 项的系数（不必求出其它各项的系数）。

\Needspace{4\baselineskip}
\textbf{解析}

\Needspace{5\baselineskip}
$(1+x)^k$ 中 $x^8$ 的系数为组合数 $C_k^8$，当 $k<8$ 时该项不存在。故所求系数为

\nopagebreak[4]
\[
C_8^8+C_9^8+C_{10}^8
\]

\Needspace{5\baselineskip}
逐项计算：$C_8^8=1$，$C_9^8=9$，$C_{10}^8=C_{10}^2=45$，故

\nopagebreak[4]
\[
C_8^8+C_9^8+C_{10}^8=1+9+45=55
\]

（也可用组合恒等式 $C_8^8+C_9^8+C_{10}^8=C_{11}^9=55$ 一步得到。）

\gkanswer{答案：$55$}

\end{answer}
\begin{answer}
\gksolhead{3}{}

\textbf{题目}：

求圆锥曲线 $2x^2-8xy-4y^2-4y+1=0$ 之焦点及准线。

\Needspace{4\baselineskip}
\textbf{解析}

\Needspace{5\baselineskip}
记 $f(x,y)=2x^2-8xy-4y^2-4y+1$，二次项矩阵

\nopagebreak[4]
\[
M=\begin{pmatrix}2&-4\\-4&-4\end{pmatrix}
\]

\textbf{求中心}：令 $f_x=4x-8y=0$，$f_y=-8x-8y-4=0$，解得中心 $C_0\left(-\dfrac13,-\dfrac16\right)$。

\textbf{判类型并转轴}：特征方程 $\det(M-\lambda I)=(2-\lambda)(-4-\lambda)-16=\lambda^2+2\lambda-24=0$，得 $\lambda_1=4$，$\lambda_2=-6$，一正一负，曲线为双曲线。对应单位特征向量：

\begin{itemize}
\item $\lambda_1=4$：$(M-4I)v=0$ 给出 $v\parallel(2,-1)$，取 $e_X=\dfrac{(2,-1)}{\sqrt5}$；
\item $\lambda_2=-6$：$(M+6I)v=0$ 给出 $v\parallel(1,2)$，取 $e_Y=\dfrac{(1,2)}{\sqrt5}$。
\end{itemize}

\Needspace{5\baselineskip}
以 $C_0$ 为新原点、$e_X$、$e_Y$ 为新轴方向，设曲线上点相对 $C_0$ 的位移为 $Xe_X+Ye_Y$。平移消去一次项（$C_0$ 是中心），转轴后二次项系数为特征值，方程化为

\nopagebreak[4]
\[
4X^2-6Y^2+\frac43=0
\]

\Needspace{5\baselineskip}
即

\nopagebreak[4]
\[
\frac{Y^2}{\;\dfrac29\;}-\frac{X^2}{\;\dfrac13\;}=1
\]

\Needspace{5\baselineskip}
实轴沿 $e_Y$ 方向，$a^2=\dfrac29$，$b^2=\dfrac13$，半焦距

\nopagebreak[4]
\[
c=\sqrt{a^2+b^2}=\sqrt{\frac29+\frac39}=\frac{\sqrt5}3
\]

\Needspace{5\baselineskip}
\textbf{焦点}：中心沿实轴方向 $\pm c\,e_Y$ 处。由 $\dfrac{\sqrt5}3\cdot\dfrac{(1,2)}{\sqrt5}=\left(\dfrac13,\dfrac23\right)$，得

\nopagebreak[4]
\[
\left(-\frac13,-\frac16\right)\pm\left(\frac13,\frac23\right)
\]

即焦点为 $\left(0,\dfrac12\right)$ 与 $\left(-\dfrac23,-\dfrac56\right)$。

\Needspace{5\baselineskip}
\textbf{准线}：双曲线 $\dfrac{Y^2}{a^2}-\dfrac{X^2}{b^2}=1$ 的准线为 $Y=\pm\dfrac{a^2}{c}$，而

\nopagebreak[4]
\[
\frac{a^2}{c}=\frac{2/9}{\sqrt5/3}=\frac{2}{3\sqrt5}
\]

\Needspace{5\baselineskip}
$Y$ 即点相对 $C_0$ 的位移在 $e_Y$ 上的投影：$Y=\dfrac{(x+\frac13)+2(y+\frac16)}{\sqrt5}$。准线方程为

\nopagebreak[4]
\[
\frac{\left(x+\frac13\right)+2\left(y+\frac16\right)}{\sqrt5}=\pm\frac{2}{3\sqrt5}
\]

\Needspace{5\baselineskip}
两边乘 $\sqrt5$：

\nopagebreak[4]
\[
\left(x+\frac13\right)+2\left(y+\frac16\right)=\pm\frac23
\]

\Needspace{5\baselineskip}
即 $x+2y+\dfrac23=\pm\dfrac23$，也就是两条直线

\nopagebreak[4]
\[
x+2y=0\qquad\text{与}\qquad x+2y+\frac43=0
\]

（验证：曲线上任一点到焦点 $\left(0,\dfrac12\right)$ 的距离与到准线 $x+2y=0$ 的距离之比为离心率 $e=\dfrac ca=\sqrt{\dfrac52}$，与准线定义一致。）

\gkanswer{答案：焦点 $\left(0,\dfrac12\right)$、$\left(-\dfrac23,-\dfrac56\right)$；准线 $x+2y=0$ 与 $x+2y+\dfrac43=0$}

\end{answer}
\begin{answer}
\gksolhead{4}{}

\textbf{题目}：

求圆锥曲线 $x^2+y^2=49$ 与 $x^2+y^2-20y+91=0$ 之公切线。

\Needspace{4\baselineskip}
\textbf{解析}

第一个圆 $x^2+y^2=49$：圆心 $O(0,0)$，半径 $r_1=7$。第二个圆 $x^2+y^2-20y+91=0$ 即 $x^2+(y-10)^2=9$：圆心 $O'(0,10)$，半径 $r_2=3$。

圆心距 $|OO'|=10$，恰等于 $r_1+r_2=7+3$，故两圆外切，切点在连心线上：$(0,7)$。外切两圆的公切线共三条：一条过切点的公切线与两条外公切线。

\Needspace{5\baselineskip}
\textbf{过切点的公切线}：过 $(0,7)$ 且垂直于连心线（$y$ 轴），即

\nopagebreak[4]
\[
y=7
\]

\Needspace{5\baselineskip}
\textbf{外公切线}：设 $y=kx+m$，记 $s=\sqrt{1+k^2}$。外公切线时两圆心在直线同侧（直线在两圆心上方），用有符号距离：

\nopagebreak[4]
\[
\frac{m}{s}=7,\qquad \frac{m-10}{s}=3
\]

\Needspace{5\baselineskip}
两式相减得 $\dfrac{10}{s}=4$，即 $s=\dfrac52$，故

\nopagebreak[4]
\[
k^2=s^2-1=\frac{25}{4}-1=\frac{21}{4},\qquad k=\pm\frac{\sqrt{21}}2
\]

\Needspace{5\baselineskip}
且 $m=7s=\dfrac{35}{2}$。两条外公切线为

\nopagebreak[4]
\[
y=\frac{\sqrt{21}}{2}x+\frac{35}{2},\qquad y=-\frac{\sqrt{21}}{2}x+\frac{35}{2}
\]

（检验：$\dfrac{35/2}{5/2}=7$，$\dfrac{35/2-10}{5/2}=3$，两圆心到直线距离分别等于各自半径。）

故所求公切线共三条：$y=7$，$y=\pm\dfrac{\sqrt{21}}2x+\dfrac{35}2$。

\gkanswer{答案：$y=7$；$y=\dfrac{\sqrt{21}}2x+\dfrac{35}2$；$y=-\dfrac{\sqrt{21}}2x+\dfrac{35}2$}

\end{answer}
\begin{answer}
\gksolhead{5}{}

\textbf{题目}：

证明：设一四面体为一平面所截，若截面是平行四边形，则此截面平行于四面体的两条棱。

\Needspace{4\baselineskip}
\textbf{解析}

设四面体为 $ABCD$，截面为平行四边形 $EFGH$，其中 $E\in AB$，$F\in BC$，$G\in CD$，$H\in DA$，四条边分别落在四个面内。要证：截面平行于一组对棱。

\Needspace{5\baselineskip}
由 $EFGH$ 是平行四边形，$EF\parallel HG$。$EF\subset$ 平面 $ABC$，而 $HG\subset$ 平面 $ACD$（$G\in CD$，$H\in DA$ 均在平面 $ACD$ 内）。$EF$ 平行于平面 $ACD$ 内的直线 $HG$，且 $EF\not\subset$ 平面 $ACD$，故

\nopagebreak[4]
\[
EF\parallel \text{平面 } ACD
\]

\Needspace{5\baselineskip}
平面 $ABC$ 与平面 $ACD$ 的交线是棱 $AC$。由线面平行的性质定理（若一条直线平行于一个平面，则过它的平面与该平面的交线平行于这条直线），

\nopagebreak[4]
\[
EF\parallel AC
\]

\Needspace{5\baselineskip}
又 $EF\subset$ 截面平面 $\pi$，且 $AC\not\subset\pi$（棱 $AC$ 不在截面内），故由线面平行的判定定理得

\nopagebreak[4]
\[
AC\parallel\pi
\]

\Needspace{5\baselineskip}
同理，由 $FG\parallel EH$，$FG\subset$ 平面 $BCD$，$EH\subset$ 平面 $ABD$，得 $FG\parallel$ 平面 $ABD$，其交线为棱 $BD$，故 $FG\parallel BD$，从而

\nopagebreak[4]
\[
BD\parallel\pi
\]

$AC$ 与 $BD$ 是四面体的一组对棱。所以平行四边形截面必平行于四面体的一组（两条）对棱。

（若截面四顶点分布在另外四条棱上，将字母相应轮换后论证完全相同。）

\gkanswer{答案：见解析版}

\end{answer}
\begin{answer}
\gksolhead{6}{}

\textbf{题目}：

设三角形 $ABC$ 中，$\angle A$ 的外角平分线与 $BC$ 边之延长线相交于 $D$，与外接圆周交于 $E$，求证：$AB\cdot AC=AE\cdot AD$。

\Needspace{4\baselineskip}
\textbf{解析}

\Needspace{5\baselineskip}
记 $\angle BAC=A$。外角平分线是过 $A$ 的直线：它与内角平分线垂直，与边 $AC$ 所成的锐角为

\nopagebreak[4]
\[
\frac{180^\circ-A}{2}=90^\circ-\frac A2
\]

该直线交 $BC$ 的延长线于 $D$，交外接圆于 $E$；射线 $AE$ 与射线 $AD$ 方向相反（$E$ 在优弧 $BAC$ 上，$D$ 在 $A$ 的另一侧）。

\Needspace{5\baselineskip}
\textbf{角度合同}：射线 $AE$ 恰平分 $AC$ 与 $AB$ 反向延长线所成的角，故 $\angle CAE=90^\circ-\dfrac A2$；且射线 $AE$ 与射线 $AB$ 分居 $AC$ 两侧，于是

\nopagebreak[4]
\[
\angle BAE=\angle BAC+\angle CAE=A+\left(90^\circ-\frac A2\right)=90^\circ+\frac A2
\]

\Needspace{5\baselineskip}
又 $AD$ 与 $AE$ 方向相反，故

\nopagebreak[4]
\[
\angle DAC=180^\circ-\angle CAE=180^\circ-\left(90^\circ-\frac A2\right)=90^\circ+\frac A2
\]

\Needspace{5\baselineskip}
因此

\nopagebreak[4]
\[
\angle BAE=\angle DAC
\]

\Needspace{5\baselineskip}
\textbf{第二个角}：$\angle AEB$ 与 $\angle ACB$ 是同弧 $AB$ 上的圆周角，故 $\angle AEB=\angle ACB$；$D$ 在 $BC$ 的延长线上，射线 $CD$ 与射线 $CB$ 同向，故 $\angle ACD=\angle ACB$。于是

\nopagebreak[4]
\[
\angle AEB=\angle ACD
\]

\Needspace{5\baselineskip}
\textbf{相似得积}：由两角相等，$\triangle ABE\sim\triangle ADC$（对应关系 $A\leftrightarrow A$，$B\leftrightarrow D$，$E\leftrightarrow C$），故

\nopagebreak[4]
\[
\frac{AB}{AD}=\frac{AE}{AC}
\]

\Needspace{5\baselineskip}
即

\nopagebreak[4]
\[
AB\cdot AC=AE\cdot AD
\]

（若 $D$ 在 $CB$ 的延长线一侧，各角度关系对称，证明相同。）

\gkanswer{答案：见解析版}

\end{answer}
\begin{answer}
\gksolhead{7}{}

\textbf{题目}：

\Needspace{5\baselineskip}
设有互相外切的二圆的半径各为 $a$ 与 $b$（$a>b$），作此二圆的两外公切线，设其夹角为 $A$，证明：

\nopagebreak[4]
\[
\sin A=\frac{4(a-b)\sqrt{ab}}{(a+b)^2}
\]

\Needspace{4\baselineskip}
\textbf{解析}

\Needspace{5\baselineskip}
设两圆圆心为 $O_1$（半径 $a$）、$O_2$（半径 $b$），$a>b$。两外切于 $T$，则连心距

\nopagebreak[4]
\[
O_1O_2=a+b
\]

\Needspace{5\baselineskip}
两条外公切线交于 $P$（外位似中心），$P$ 在连心线 $O_1O_2$ 上。设一条公切线分别切两圆于 $M_1$、$M_2$，则 $O_1M_1\perp M_1M_2$，$O_2M_2\perp M_1M_2$，故 $O_1M_1\parallel O_2M_2$，直角三角形 $\triangle PO_1M_1\sim\triangle PO_2M_2$：

\nopagebreak[4]
\[
\frac{PO_1}{PO_2}=\frac{O_1M_1}{O_2M_2}=\frac ab
\]

\Needspace{5\baselineskip}
又 $PO_1-PO_2=O_1O_2=a+b$，联立解得

\nopagebreak[4]
\[
PO_1=\frac{a(a+b)}{a-b},\qquad PO_2=\frac{b(a+b)}{a-b}
\]

\Needspace{5\baselineskip}
由对称性，连心线 $PO_1$ 平分两条外公切线的夹角 $A$，故切线与连心线的夹角为 $\dfrac A2$。在 $\mathrm{Rt}\triangle PO_1M_1$ 中，

\nopagebreak[4]
\[
\sin\frac A2=\frac{O_1M_1}{PO_1}=\frac{a}{\dfrac{a(a+b)}{a-b}}=\frac{a-b}{a+b}
\]

\Needspace{5\baselineskip}
于是

\nopagebreak[4]
\[
\cos\frac A2=\sqrt{1-\sin^2\frac A2}=\frac{\sqrt{(a+b)^2-(a-b)^2}}{a+b}=\frac{\sqrt{4ab}}{a+b}=\frac{2\sqrt{ab}}{a+b}
\]

\Needspace{5\baselineskip}
由二倍角公式

\nopagebreak[4]
\[
\sin A=2\sin\frac A2\cos\frac A2=2\cdot\frac{a-b}{a+b}\cdot\frac{2\sqrt{ab}}{a+b}
\]

\Needspace{5\baselineskip}
即

\nopagebreak[4]
\[
\sin A=\frac{4(a-b)\sqrt{ab}}{(a+b)^2}
\]

\gkanswer{答案：见解析版}

\end{answer}
\begin{answer}
\gksolhead{8}{}

\textbf{题目}：

\Needspace{5\baselineskip}
在三角形 $ABC$ 中，$r$ 与 $R$ 分别为其内切圆与外接圆的半径，求证：

\nopagebreak[4]
\[
r=4R\sin\frac A2\sin\frac B2\sin\frac C2
\]

\Needspace{4\baselineskip}
\textbf{解析}

记三角形三边 $a=BC$，$b=CA$，$c=AB$，半周长 $s=\dfrac{a+b+c}2$，面积 $\Delta$，内切圆半径 $r$，外接圆半径 $R$。

\Needspace{5\baselineskip}
\textbf{半角公式}：

\nopagebreak[4]
\[
\sin\frac A2=\sqrt{\frac{(s-b)(s-c)}{bc}},\qquad
\sin\frac B2=\sqrt{\frac{(s-c)(s-a)}{ca}},\qquad
\sin\frac C2=\sqrt{\frac{(s-a)(s-b)}{ab}}
\]

\Needspace{5\baselineskip}
\textbf{三式相乘}：分子恰为 $(s-a)(s-b)(s-c)$，分母为 $abc\,s$ 的幂次——分母中 $bc\cdot ca\cdot ab=a^2b^2c^2$，开方后得

\nopagebreak[4]
\[
\sin\frac A2\sin\frac B2\sin\frac C2=\frac{(s-a)(s-b)(s-c)}{abc}
\]

\Needspace{5\baselineskip}
\textbf{用面积公式代换}：由海伦公式

\nopagebreak[4]
\[
\Delta^2=s(s-a)(s-b)(s-c)\ \Rightarrow\ (s-a)(s-b)(s-c)=\frac{\Delta^2}{s}
\]

\Needspace{5\baselineskip}
又由三角形面积的外接圆表达式与内切圆表达式

\nopagebreak[4]
\[
abc=4R\Delta,\qquad \Delta=rs
\]

\Needspace{5\baselineskip}
代入：

\nopagebreak[4]
\[
\sin\frac A2\sin\frac B2\sin\frac C2=\frac{\Delta^2/s}{4R\Delta}=\frac{\Delta}{4Rs}=\frac{rs}{4Rs}=\frac r{4R}
\]

\Needspace{5\baselineskip}
整理即得

\nopagebreak[4]
\[
r=4R\sin\frac A2\sin\frac B2\sin\frac C2
\]

\gkanswer{答案：见解析版}

\end{answer}